% Literature Review revised for the multiview RGB-D / retargeting / 3DGS topic.
% Uses the original thesis citation keys, plus mimicgen, dexretargeting,
% robowheel and opencvcharuco from the previously supplied references.bib.
% Only dexteleop0 is deliberately not cited: its supplied entry identifies
% no author, paper, venue, DOI or URL and says "Citation to be confirmed".
% Bibliographic corrections to check in the existing .bib (not new claims):
% - dexflow: Lin et al., "DexFlow: A Unified Approach for Dexterous Hand
%   Pose Retargeting and Interaction", arXiv:2505.01083 (2025).
% - keyobj: Xin, Yu, Jiang, Zhang and Li, arXiv:2506.09384 (2025).
% - gaussianfactory: arXiv:2609.21112 is a 2026 paper, not a 2024 paper.
% - legs: arXiv:2606.01458 is a 2026 paper, not a 2024 paper.
% - robosimgs: arXiv:2510.10637 is a 2025 paper, not a 2024 paper.
% - alihandeye: the supplied title identifies Enebuse et al. (2022),
%   PLOS ONE 17(10):e0273261, DOI 10.1371/journal.pone.0273261.
% - hrbf: Xu, Nan, Zhou, Wang and Wang (2022), ACM TOG 41(3), article 35,
%   DOI 10.1145/3516521; the supplied author/year fields are incomplete.
% - ridgerun: use the precise RR Syncer technical documentation discussed here:
%   https://developer.ridgerun.com/wiki/index.php/GStreamer_Buffer_Synchronization/Documentation/GStreamer
% Existing section labels are preserved for compatibility with the thesis.

\chapter{Literature Review}
\label{ch:lit}

\section{Preamble}
\label{sec:lit-preamble}

Generating robot demonstrations from human recordings connects several
research problems: recovering hand--object interaction, adapting that interaction
to another embodiment, reconstructing the scene, and producing observations
consistent with the resulting robot motion. The literature contains both
specialised methods for these problems and integrated data-generation systems.
The purpose of this review is to identify the representations and assumptions
that can support the proposed ViKi workflow, and the measurements needed to
assess their combination.

The chapter retains work on teleoperation, calibration, geometric fusion, and
downstream learning because these areas establish requirements on the generated
data even when they are not themselves the final application of the thesis.
Teleoperation offers reference retargeting formulations; calibration and fusion
determine the metric consistency of the source; learning systems define the
observations and action labels that a usable dataset must supply. Their
relevance is therefore methodological as well as comparative.

Section~\ref{sec:lit-hand} examines human-motion representations and acquisition
interfaces. Section~\ref{sec:lit-retarget} reviews retargeting objectives,
temporal processing, and solvers. Section~\ref{sec:lit-3dgs} discusses 3DGS-based
rendering and demonstration generation. Sections~\ref{sec:lit-repr}
and~\ref{sec:lit-calib} address task coordinates and capture geometry.
Section~\ref{sec:lit-dataset} considers dataset construction and assessment.
Section~\ref{sec:positioning} positions ViKi relative to the closest systems,
and Section~\ref{sec:lit-conclusion} synthesises the design implications.
The review is a narrative comparison of the cited publications and technical
documentation, rather than evidence of an exhaustive absence of related work.

\section{Representing the Human Hand}
\label{sec:lit-hand}

\subsection{Articulated hand and body reconstruction}

The representation of a demonstration determines which aspects of manipulation
can be transferred. A wrist trajectory captures global hand motion, whereas
finger articulation and hand--object geometry are needed to distinguish
contacts that share a similar wrist pose. DexMV combines hand and object pose
estimation with demonstration translation for dexterous robot learning
\cite{dexmv}. OKAMI uses a single RGB-D video to derive a manipulation plan,
retargeting body and hand motion with object-aware adaptation
\cite{okami}. These works establish that human reconstruction and robot action
generation are connected stages, rather than interchangeable outputs.

DexCanvas provides an instrumented reference for richer interaction capture
\cite{dexcanvas}. Its acquisition system combines 22 infrared motion-capture
cameras, two RGB-D cameras, and markers on the hand and objects. The observations
are fitted to MANO and used in a real-to-simulation process that reconstructs
contact-related forces through reinforcement learning. Its relevance to ViKi
is twofold: articulated representations retain information needed by dexterous
end-effectors, and contact annotations require evidence beyond a visually
plausible hand trajectory. Its marker-based acquisition also differs from the
proposed markerless reconstruction route.

A fitted hand model can regularise noisy or partially observed landmarks, but
it introduces assumptions about shape, articulation, and correspondence.
Landmark error, mesh-fit error, and the accuracy of inferred contacts should
therefore be evaluated separately. A parametric model used to create reference
annotations need not be the model used by the evaluated perception pipeline.
For ViKi, a useful representation preserves the available articulated
observations while allowing target-specific reductions for simpler grippers.
This avoids committing the entire source dataset to a single end-effector.

\subsection{Compact arm and gripper representations}

For a parallel-jaw target, a pose in $\SE$ together with an opening command
provides a compact action description. It can represent approach, transport,
and approximate closing and release timing. It does not encode all finger
contacts, in-hand manipulation, or the forces exerted by the demonstrator.
Consequently, a scalar pinch signal can be an appropriate target abstraction
without being a lossless representation of the original hand motion.

AnyTeleop supports different arms, hands, and camera configurations within a
vision-based teleoperation framework \cite{anyteleop}. Open-TeleVision maps
operator arm and hand motion to robot embodiments while providing immersive
visual feedback \cite{opentv}. Bunny-VisionPro combines bimanual dexterous
teleoperation with haptic feedback and collision and singularity avoidance
\cite{bunny}. These systems motivate separating arm motion, finger
correspondence, and embodiment-specific constraints instead of treating
retargeting as a uniform rescaling of a human skeleton.

ACE further demonstrates that one collection interface can serve humanoid
hands, arm--hand, arm--gripper, and quadruped--gripper configurations
\cite{ace}. It combines a hand-facing camera with an exoskeleton to observe
finger and wrist motion. Thus, teleoperation does not necessarily require a
completely different operator interface for every robot. The distinction
relevant to ViKi is whether a source demonstration can be recorded independently
of a target robot and subsequently converted into several robot datasets.

The retargeting study of Xin et al. analyses the contribution of alternative
hand-configuration objectives through ablations and manipulation experiments
\cite{keyobj}. Its results support evaluating the choice and weighting of
correspondences, rather than assuming that maximal similarity to every human
joint is always desirable. Conclusions from a dexterous-hand ablation do not,
by themselves, determine the correct wrist weight for a different arm--gripper
system. That weight depends on which task relationships the target must retain.

\subsection{Egocentric capture and human--robot co-training}

EgoMimic records egocentric human observations and 3D hand motion using Project
Aria glasses and co-trains on human and robot demonstrations
\cite{egomimic}. DexWild likewise uses human hand interactions collected across
diverse environments, together with robot data, to improve policy transfer
\cite{dexwild}. These are relevant alternatives to converting every human
frame into a photorealistic image of a replacement robot. They show that human
data can be useful through shared representations and co-training, although
that route has different dataset requirements from explicit robot-action and
robot-view synthesis.

Vision-Based Hand Shadowing provides a closer analytical video-to-robot
workflow \cite{shadowing}. It uses an egocentric RGB-D camera, MediaPipe
landmarks, depth lifting, damped-least-squares inverse kinematics, and a
thumb--index-based gripper mapping for the SO-ARM101. Its reported degradation
under occlusion makes visibility a relevant evaluation condition, not merely
a reason to prefer more cameras. Multiview capture can add geometric evidence,
but any improvement must be measured with controlled calibration, timing, and
view selection.

\section{Retargeting Formulations}
\label{sec:lit-retarget}

\subsection{Correspondences, objectives, and confidence}

Retargeting converts an observed human configuration or interaction into robot
variables subject to a different kinematic structure. Position-based methods
match selected points; vector-based methods match relative directions or
lengths; task-space methods track a wrist, tool, or object relationship.
AnyTeleop and the accompanying \texttt{dex-retargeting} toolkit illustrate
several such formulations \cite{anyteleop, dexretargeting}. Their common
principle is to select geometric correspondences that the target can realise,
while accounting for joint limits and temporal behaviour.

These objectives can be combined, but their weights have different meanings.
A task weight expresses the importance of a relationship, while a perception
weight expresses the reliability of its observation. A contact-dependent
weight need not be a calibrated uncertainty estimate. In ViKi, this distinction
matters when an occluded fingertip remains important to the task: increasing
its task importance cannot compensate for a missing measurement. Confidence
weighting therefore needs an evaluation of its own effects on tracking,
missing intervals, and failure cases.

The choice between wrist, fingertip, and object-relative residuals also depends
on the target end-effector. For a parallel gripper, aperture and tool pose may
be sufficient for a restricted pick-and-place task. For an Allegro-type hand,
additional finger correspondences and grasp objectives are required. The
representation should expose these differences rather than imply that one
pinch mapping solves cross-morphology transfer in general.

\subsection{Temporal objectives and sequential processing}

Offline access to a complete demonstration permits temporal filtering,
sliding-window optimisation, or optimisation over a complete robot trajectory.
A sequential alternative reconstructs and smooths task-space motion, solves
inverse kinematics frame by frame, and optionally filters joint commands.
Because inverse kinematics is nonlinear and may switch solution branches,
smooth Cartesian targets need not imply smooth joint motion. Conversely,
filtering joint commands can change the end-effector path or invalidate
previously checked clearances.

DexFlow combines hand-pose alignment with temporal consistency and
hand--object interaction refinement \cite{dexflow}. Its formulation includes
sliding-window processing and motion prediction, and the authors discuss
higher-order continuity in connection with regularisation. This supports
examining temporal terms alongside pose alignment. It does not establish that
any implementation using a smoothness penalty automatically satisfies a
velocity, acceleration, or jerk bound.

A useful distinction is between a penalty and a constraint. Penalising large
joint differences expresses a preference; enforcing a limit requires an
explicit bound or an acceptance test on the resulting trajectory. Similarly,
a robust residual may reduce the influence of an outlier but does not identify
an unobserved contact or recover information lost during occlusion. The
comparison of preprocessing and joint-space regularisation should therefore
report both geometric tracking error and the motion derivatives affected by
the chosen method.

\subsection{Offline transfer and interaction-aware alternatives}

HOWTransfer estimates temporally consistent hand motion, localises contact
intervals, converts grasp intent into parallel-jaw grasp hypotheses, and edits
the resulting trajectories \cite{howtransfer}. Its hand-centric formulation
is relevant when reliable contact cues are available without complete explicit
object-state tracking. It offers a distinct comparison to a pipeline that
first reconstructs object poses and then imposes object-relative constraints.

ORION instead derives an object-centric manipulation plan from a human RGB-D
video \cite{orion}. RoboWheel integrates human hand--object reconstruction,
physical-plausibility refinement, cross-embodiment retargeting, and simulation
augmentation \cite{robowheel}. These systems demonstrate that offline
reconstruction, transfer, and data generation have already been studied in
combination. The relevant comparison is between their assumptions and
representations, not a division into integrated ViKi and otherwise isolated
prior components.

For ViKi, hand-only, object-relative, and contact-aware retargeting can be
compared using the same source motion and robot geometry. Contact-aware
checking must distinguish permitted gripper--object contact from unintended
penetration. It must also state which scene surfaces enter the collision
model. A self-collision-free arm trajectory is not necessarily clear of the
table, and imposing a recorded object trajectory does not show that the robot
would physically cause that motion.

\subsection{Inverse-kinematics libraries and trajectory optimisation}

Pink provides a task-based differential inverse-kinematics interface built on
Pinocchio \cite{pink}. Weighted tasks and supported limits or barriers can be
assembled into local quadratic programs. This is useful for expressing
end-effector targets and kinematic restrictions. However, a local program at
one timestep and a batch problem coupling many timesteps have different
variables and numerical properties. Choosing Pink does not, on its own,
implement a complete trajectory-level objective.

The solver description should therefore distinguish local inverse kinematics,
temporal coupling, and final verification. Joint limits, velocity limits,
collision geometry, and numerical tolerances are properties of the configured
application as well as the library. The resulting trajectory needs checks in
the original geometry, particularly when linearised collision constraints or
post-processing are used. This motivates retargeting-feasibility measurements
without equating solver convergence with successful manipulation.

\section{3DGS for Robot Learning}
\label{sec:lit-3dgs}

3DGS represents scene appearance using explicit Gaussian
primitives with optimised position, covariance, opacity, and view-dependent
colour. Kerbl et al. demonstrate high-quality novel-view synthesis with an
efficient splatting renderer \cite{gaussiansplatting}. Its explicit spatial
representation supports component editing and compositing, making it useful
for generating observations beyond the original camera views. Robot actions,
contact geometry, and time-varying object states must still come from another
part of the system.

\subsection{Augmenting recorded robot demonstrations}

RoboSplat starts from a robot demonstration and multiview scene images, then
edits reconstructed Gaussian components to generate additional training data
\cite{robosplat}. Its six augmentation dimensions are object pose, object
type, camera view, scene appearance, lighting, and embodiment. Thus, it already
addresses embodiment variation as well as viewpoint and appearance changes.
The distinction from the proposed ViKi acquisition stage is that its source
motion comes from robot teleoperation. Its geometric editing operations should
not be described as a general diffusion-based trajectory generator.

This example separates two forms of expansion. Re-rendering an existing
trajectory changes the observation distribution; adapting object poses or
robot embodiment may also require changing the action trajectory. Both can
increase dataset coverage, but they create different validity requirements.
A rendered robot must agree with its exported state, and a new geometric
configuration must be checked for reachability and interaction consistency.

\subsection{Simulation-based data generation and hybrid rendering}

SplatSim uses Gaussian splatting for photorealistic observations coupled to a
simulator \cite{splatsim}. Its experiments include demonstrations collected
through simulated teleoperation or generated by motion planners, followed by
policy training and real-world deployment. It is therefore relevant both as a
rendering approach and as a synthetic-data workflow. It does not depend on
recovering motion from a passive human recording.

RoboSimGS combines Gaussian scene appearance with interactive mesh assets and
physics simulation \cite{robosimgs}. Its pipeline includes inference of
object properties and articulated structures using a multimodal model. The
method is relevant to ViKi's separation of appearance and interaction
geometry: photorealistic background reconstruction and usable collision assets
are related but different reconstruction products.

LEGS uses a mesh foreground for the robot and objects over a reconstructed
3DGS background, together with procedural motion primitives for humanoid
loco-manipulation \cite{legs}. It also studies alignment between rendered
appearance and the deployment camera. Its separation of motion and appearance
allows demonstrations to be re-rendered under changed scene conditions. The
human-motion source differs from ViKi's proposal, but the hybrid rendering
architecture and appearance calibration remain directly relevant.

These systems establish that a 3DGS-based pipeline need not render every
component as Gaussians. A mesh foreground can supply articulation and depth
while a splat background supplies appearance. Such a combination still requires
metric alignment, compatible camera models, and correct visibility ordering.
A visually detailed background does not compensate for a misplaced robot or
an incorrect object state.

\subsection{Planning demonstrations from reconstructed scenes}

GaussianFactory generates demonstrations from a video scan by reconstructing
editable 3DGS assets and planning task motion geometrically
\cite{gaussianfactory}. Grasp formation uses a pretrained contact model, and
the work includes downstream policy evaluation. Its input is scene structure
rather than a human manipulation trajectory. This offers an alternative
motion source for dataset generation, not evidence that procedurally generated
motion is inherently inferior. A comparison with demonstration transfer should
consider the tasks covered, preparation effort, and the interaction information
that must be specified or inferred in each case.

\subsection{Human demonstrations and novel-view synthesis}

Tool-as-Interface is a particularly relevant precedent \cite{toolif}. It uses
two RGB views, MASt3R reconstruction, and Gaussian splatting to augment human
tool-use observations across viewpoints. Tool poses provide task-space action
labels, while segmentation removes embodiment-specific regions from human
and robot observations. This combines human demonstration reuse with 3D-aware
visual synthesis without requiring the complete human hand to be reconstructed.
Its shared-tool representation differs from explicitly rendering target robot
embodiments, but its existence rules out treating human-to-robot transfer and
Gaussian view augmentation as unrelated lines of work.

Real2Render2Real uses a smartphone object scan and a monocular human
manipulation video to recover object-part motion and generate robot trajectories
\cite{real2render2real}. It uses 3DGS for asset construction and tracking,
then converts reconstructed assets to textured meshes for rendering in
IsaacLab with dynamics disabled. The project supports retargeting across robot
embodiments and variation in camera and scene conditions. Its sensing and final
rendering differ from ViKi's proposed multiview RGB-D route, but its reuse of
human-demonstrated object motion is a close methodological precedent.

RoboPaint more directly combines multiview RGB-D human capture, retargeting,
and 3DGS-based visual synthesis \cite{robopaint}. It records glove kinematics
and tactile signals alongside camera observations, uses wrist markers, and
retargets motion with kinematic and tactile objectives. Its hybrid renderer
combines a static 3DGS background with mesh robots and moving objects, supporting
novel views and scene edits. The principal acquisition distinction proposed
for ViKi is recovery from camera observations without worn hand markers or
instrumented gloves. That distinction is a design choice to evaluate, rather
than proof that the broad integration is new.

\subsection{Reconstruction boundaries and evaluation}

The reviewed methods use direct splat rendering, hybrid rendering, or
3DGS-to-mesh conversion for different purposes. These choices affect how
objects are segmented, how trajectories drive the scene, and which surfaces
are observable from new cameras. A manipulation recording also contains a
moving demonstrator and possibly moving objects; a static reconstruction
cannot simply be assumed to reproduce all of them consistently over time.

ViKi therefore needs an explicit decomposition into the background, moving
object assets, and replacement robot, together with a method for removing the
human from generated observations. Reconstruction quality can be measured
against held-out views of the same scene state. Replacement-robot compositing
requires additional checks on depth, boundaries, temporal consistency, and
agreement with exported states. A human image is not a paired pixel-level
reference for an image containing a different robot. These distinctions keep
visual fidelity separate from the correctness of the generated action labels.

\section{Task Representation}
\label{sec:lit-repr}

\subsection{World, object, and tool coordinates}

A trajectory expressed in the capture workspace is straightforward to inspect,
but replaying it after an object has moved may miss the intended interaction.
Object-centric and tool-centric representations instead preserve selected
relations between the end-effector and the task. ORION extracts an
object-centric plan, and SPOT represents task motion using an object-pose
trajectory relative to a target \cite{orion, spot}. Their relevance is the
separation of task structure from the particular joint trajectory used to
realise it.

OKAMI occupies a related position: object-aware adaptation modifies recovered
human motion for new object locations \cite{okami}. Vi-RoIL expresses wrist
trajectories and key-frame hand poses in object coordinates after monocular
hand--object reconstruction with a wrist-mounted marker cube
\cite{viroil}. Its reported qualitative simulation results support discussion
of coordinate consistency, while providing a different evidence level from a
quantitative real-robot manipulation benchmark.

Tool-as-Interface uses the tool as a shared physical interface between the
human and robot \cite{toolif}. This is particularly useful when the tool can
be rigidly attached or grasped with a known transform. It also illustrates the
limits of a universal wrist representation: task-relevant motion may be more
naturally specified at the tool than at an anatomical hand joint.

\subsection{Flow representations and non-rigid interaction}

NovaFlow generates task videos from descriptions, extracts actionable object
flow, and uses that flow for robot planning \cite{novaflow}. For rigid objects,
it derives relative poses; for deformable objects, it uses flow as an objective
for model-based planning. Its motion source is generated video rather than a
recorded human demonstration. Nevertheless, it is relevant to the representation
choice because it shows how a flow-based interface can extend beyond a single
rigid-pose trajectory.

A rigid-object ViKi pipeline should therefore identify its object assumptions
without implying that all demonstration-driven transfer requires rigidity.
Cloth, fluids, and changing topology require a different motion and interaction
model. Even for rigid objects, symmetry and occlusion can make orientation
ambiguous. A stable-looking object track may still impose an incorrect
relative grasp when used downstream.

\subsection{Implications for reusable demonstrations}

Object-relative motion can be stored alongside world-frame observations, as
considered in Section~\ref{sec:objcentric}. Retaining both allows comparison
of direct hand tracking with interaction-aware retargeting and preserves the
original evidence when an object estimate fails. Object-relative coordinates
do not create invariance to arbitrary scene changes: the new layout must still
admit the interaction and a feasible robot path.

The relevant preservation target is task-dependent hand--object motion and
interaction timing. A robot with a different reach or end-effector may need a
different approach or grasp. Consequently, cross-embodiment reuse should report
which relationships were preserved, which were adapted, and which source
episodes could not be transferred under the declared constraints.

\section{Multiview RGB-D Capture Geometry}
\label{sec:lit-calib}

\subsection{Reference data and camera configuration}

HO-Cap provides eight RealSense views and a HoloLens recording, together with
semi-automatic hand and object annotations \cite{hocap}. It is relevant both
as a source of manipulation recordings and as an external reference for
perception assessment. The annotations are not an independent optical-motion-
capture measurement of ViKi's own rig. Dataset-specific performance and local
capture accuracy remain separate quantities.

Camera count, baseline, overlap, viewing direction, and timing jointly affect
what can be reconstructed. Additional views may reveal occluded fingers or
object surfaces, but inconsistent calibration can also introduce conflicting
constraints. A camera subset selected on the evaluation frames measures a
development choice rather than held-out generalisation. This motivates
controlled comparisons of camera configurations without making two cameras
a defining property of the proposed method.

\subsection{Fiducials and multiview registration}

ChArUco combines marker identification with chessboard-corner localisation and
supports calibration from partially visible boards \cite{opencvcharuco}.
The appropriate target depends on visibility, workspace size, and whether
intrinsics, extrinsics, or a workspace anchor are required. A small image
reprojection error is one useful diagnostic; it does not alone establish
metric agreement between depth clouds or accurate camera-to-robot registration.

Romeo et al. compare two- and three-dimensional calibration procedures for
Azure Kinect data \cite{romeo}. Their reported correspondence errors differ
between colour and infrared reconstructions and between static and dynamic
experiments. Darwish et al. calibrate a five-Azure-Kinect system with a
ChArUco-coded cube and combine reprojection with point-to-plane and
patch-to-plane terms \cite{darwish}. Both studies motivate evaluating spatial
consistency in the reconstructed geometry, rather than treating an image-space
residual as the only calibration criterion.

Kim et al. initialise multiview registration using body tracking and refine it
through feature matching and mismatch rejection \cite{kimcal}. Their feature
inlier ratios describe the refinement procedure, not directly the absolute
accuracy of a manipulation trajectory. Bu and Lee similarly combine
body-tracking-based coarse alignment with ICP refinement \cite{bulee}.
Such approaches offer markerless calibration alternatives, but depend on
initial alignment and sufficient geometric information for refinement. They
remain relevant to camera placement even if ViKi uses a fiducial workspace
anchor.

\subsection{Sensor characterisation and metric accuracy}

T\"olgyessy et al. evaluate Azure Kinect precision, systematic error, warm-up,
material response, and depth artifacts \cite{tolgyessy}. Kurillo et al.
compare Azure Kinect and Kinect v2 using reference geometry and analyse
variation with range and image position \cite{kurillo}. These studies show
why depth quality should be characterised under the actual sensing conditions:
central planar surfaces, object boundaries, reflective materials, and partially
occluded hands do not present the same estimation problem.

The RealSense D400 documentation likewise distinguishes depth accuracy, fill
rate, spatial noise, and temporal noise under stated test conditions
\cite{intel}. These are useful specifications for choosing and testing capture
settings, but they are not interchangeable with multiview registration error,
3D hand-joint error, or robot TCP error. Nor does an upper specification bound
establish a lower error floor for all measurements. ViKi's metric claims need
controls on the relevant rig and clearly defined reference quantities.

\subsection{Depth lifting, triangulation, and fusion}

Depth lifting estimates a surface location from an image coordinate and a depth
sample. Multiview triangulation instead combines viewing rays, optionally with
depth constraints. For hand tracking, the visible skin surface and anatomical
joint centre are not identical. For scene reconstruction, missing depth and
mixed pixels at boundaries can affect geometry and subsequent occlusion
rendering. These distinctions motivate separating sparse landmark fusion from
dense surface reconstruction.

The ray-casting fusion described in U.S. Patent 10,818,071 aggregates
signed-distance information from several depth images and includes weighting
related to view geometry and visibility \cite{fusionpatent}. It is relevant
as a concrete geometric fusion design, rather than comparative evidence that
a particular weighting law is optimal for hand landmarks.

Zeng et al. learn confidence-dependent RGB-D feature fusion for surface-normal
estimation \cite{zeng}. Their contribution concerns weighting depth features
before combining them with colour features in a single-image prediction
network. It supports the broader treatment of depth reliability, but is not
a direct validation of multi-camera joint averaging. HRBF-Fusion uses
continuous implicit surfaces for camera tracking and derives curvature and
confidence information for RGB-D fusion \cite{hrbf}. It provides a geometric
alternative concerned with reconstruction quality and tracking robustness.

These approaches justify considering uncertainty and local geometry when
combining observations. For ViKi, the reliability signal still needs to be
defined and tested: detector confidence, depth validity, reprojection agreement,
and interpolated samples represent different kinds of evidence. Filling a
trajectory gap increases completeness without adding an observation, so
coverage should be reported alongside reconstruction error.

\subsection{Exposure synchronisation and software alignment}

Azure Kinect supports wired multi-device synchronisation and configurable
exposure delays; its documentation also addresses interference between depth
emitters \cite{mskinect}. RealSense documentation describes its own capture
and synchronisation interfaces \cite{intel}. The applicable trigger modes and
timestamp domains must be established for the actual device combination,
rather than inferred from the fact that both families expose timestamps.

RidgeRun's RR Syncer documentation illustrates a separate software layer:
aligning video buffers by presentation timestamps, applying a synchronisation
window, and dropping unmatched buffers \cite{ridgerun}. Buffer alignment is
useful for managing streams, but does not by itself establish simultaneous
exposure. Similarly, a host monotonic clock measures arrival or application
timing unless a mapping to sensor acquisition time has been established.

For heterogeneous capture, the defensible requirement is a stated
synchronisation procedure with measured residual offsets and dropped-frame
behaviour. The literature does not justify declaring hardware coordination
between camera families impossible or treating host-side nearest-timestamp
pairing as the only available mechanism. Timing accuracy also matters to
3DGS reconstruction when moving hands and objects are observed from several
cameras.

\subsection{Camera-to-robot registration}

The generated robot must share a metric frame with the reconstructed scene.
A workspace anchor can connect fixed cameras to a calibrated robot-base pose;
hand--eye calibration is another relevant tool when suitable camera and robot
motions are available. The comparative study by Enebuse et al. finds that
hand--eye algorithm behaviour depends on translation and rotation noise and
on the motion range used for calibration \cite{alihandeye}. Its relevance
is the need to match and test the procedure under the intended conditions,
rather than to declare one solver universally preferable.

For dataset generation without a physical target robot, the chosen base pose
is a virtual placement whose transform and reachability still need to be
recorded. For later deployment, an independently established real camera-to-
robot transform is required. Neither case is validated solely by a plausible
overlay of the rendered robot on the captured scene.

\section{Dataset Generation and Downstream Requirements}
\label{sec:lit-dataset}

\subsection{Demonstration reuse and episode structure}

MimicGen adapts segments of a limited set of robot demonstrations to generate
additional demonstrations under new scene configurations \cite{mimicgen}.
Together with the human-source pipelines reviewed above, it establishes
reuse as a general data-generation strategy. The important distinctions are
the source of task motion, the assumptions used for adaptation, and the checks
applied to the generated episodes.

For ViKi, a learning episode needs temporally consistent observations, robot
states, and actions, with explicit units and coordinate conventions. The
rendered camera pose and scene state must correspond to the exported motion.
Changing the robot can change action dimension and limits; changing a camera
changes observations even when the joint trajectory is unchanged. A trajectory
bundle and a rendered observation--action dataset are therefore distinct
outputs. Compatibility with a format such as LeRobot concerns organisation
and interfaces, not certification of physical task success.

\subsection{Learning architectures as downstream consumers}

ACT, Diffusion Policy, and vision-language-action models motivate retaining
synchronised observations and action sequences, even though training them is
outside the present evaluation. OpenVLA-OFT studies action decoding,
continuous action representations, action chunking, and fine-tuning objectives
\cite{openvlaoft}. Its relevance is that action representation and temporal
organisation affect how a dataset can be consumed. Its policy results are not
evidence that a new synthetic dataset will be useful without a corresponding
learning experiment.

This downstream perspective also favours retaining robot identity, camera
calibration, source episode, retargeting parameters, and editing provenance.
Those fields allow consumers to distinguish new observations of an existing
action sequence from genuinely adapted motion, and make failures traceable to
the relevant pipeline stage.

\subsection{Curation and evidence of quality}

Demo-SCORE filters heterogeneous demonstration data using a classifier trained
and cross-validated on successful and unsuccessful policy rollouts
\cite{demoscore}. It is relevant because increasing demonstration count can
retain unreliable strategies as well as useful variation. Its quality signal,
however, depends on online experience. A pipeline without such rollouts cannot
claim equivalent success-based screening from trajectory smoothness alone.

ViKi can assess reconstruction error, geometric feasibility, rendering
consistency, and coverage before policy training. These are useful curation
criteria with a narrower interpretation than demonstrated manipulation
success. Geometric acceptance should specify which constraints were tested,
and rejected episodes should remain visible in the generation statistics.
A visually plausible sequence or a valid file schema cannot substitute for
that information.

\subsection{Factorial expansion, diversity, and generation effort}

With $M$ embodiments, $K$ camera viewpoints, and $P$ appearance configurations,
a source episode can define up to $M K P$ candidate configurations. The actual
yield depends on reachability, grasp compatibility, scene observability, and
the stated acceptance checks. Multiple views of one motion improve observation
coverage without becoming independent manipulation strategies. Appearance
changes may increase visual diversity without changing task geometry.

Consequently, dataset characterisation should report requested and accepted
combinations, failure categories, camera coverage, trajectory variation, and
processing time. Feature-space distances can provide supplementary measures
of visual variation when the encoder and sampling procedure are fixed. They
do not alone establish learning benefit or coverage of additional skills.
Likewise, a cost comparison with teleoperation needs measured operator,
calibration, annotation, scanning, computation, and rejection costs. Savings
cannot be inferred from the nominal expansion factor alone.

\section{Positioning of the Proposed System}
\label{sec:positioning}

The closest comparisons depend on the stage under consideration. Instrumented
capture systems help assess what interaction information can be recovered;
teleoperation systems supply retargeting alternatives; geometric reconstruction
and rendering systems define visual-generation choices. Table~\ref{tab:positioning}
compares representative integrated workflows by their motion source and visual
representation. It avoids equating a planned ViKi capability with a capability
already evaluated by another system.

\begingroup
\small
\setstretch{1.15}
\setlength{\tabcolsep}{4pt}
\begin{longtable}{@{}>{\raggedright\arraybackslash}p{3.2cm}>{\raggedright\arraybackslash}p{4.9cm}>{\raggedright\arraybackslash}p{6.7cm}@{}}
\caption{Motion Sources, Representations, and Visual Generation in Related Work}
\label{tab:positioning}\\
\toprule
System & Motion source and representation & Visual generation and relation to ViKi \\
\midrule
\endfirsthead
\toprule
System & Motion source and representation & Visual generation and relation to ViKi \\
\midrule
\endhead
\bottomrule
\endfoot
HOWTransfer \cite{howtransfer} & Human hand motion and contact intervals; grasp hypotheses and trajectory editing. & Relevant offline transfer alternative; not the same objective as synthesising robot-view datasets. \\[3pt]
Vi-RoIL \cite{viroil} & Monocular human motion with marker-assisted calibration; object-relative representation. & Relevant geometric front-end; reported evaluation is qualitative simulation. \\[3pt]
Tool-as-Interface \cite{toolif} & Two-view human RGB recordings; tool-centric actions. & Gaussian novel-view augmentation and embodiment masking, rather than replacement-robot rendering. \\[3pt]
Real2Render2Real \cite{real2render2real} & Human video and object-part motion; robot-specific grasping and IK. & 3DGS-to-mesh assets and synthetic rendering, including embodiment reuse. \\[3pt]
RoboPaint \cite{robopaint} & Multiview RGB-D, glove kinematics, tactile data, and wrist markers. & Human-to-robot retargeting with a 3DGS background and mesh foreground. \\[3pt]
RoboWheel \cite{robowheel} & Human interaction reconstruction and cross-embodiment retargeting. & Simulation augmentation; integrated precedent for human-source data generation. \\[3pt]
RoboSplat \cite{robosplat} & Robot demonstration with geometric augmentation. & Editable Gaussians for embodiment, viewpoint, and scene variation. \\[3pt]
SplatSim \cite{splatsim} & Simulated teleoperation or planned demonstrations. & Gaussian observations coupled to simulation; includes policy-transfer evaluation. \\[3pt]
GaussianFactory \cite{gaussianfactory} & Grasp and trajectory planning from a scene scan. & Editable 3DGS scene; planning supplies the motion. \\[3pt]
LEGS / RoboSimGS \cite{legs, robosimgs} & Procedural motion or simulated interaction. & Hybrid Gaussian appearance and mesh foreground; different motion-generation assumptions. \\[3pt]
ViKi (proposed) & Markerless multiview RGB-D hand--object motion and embodiment-specific retargeting. & Proposed robot-view rendering and controlled expansion, assessed through stage and integration checks. \\
\end{longtable}
\endgroup

RoboPaint already combines the broad components in the original motivation,
while Tool-as-Interface and Real2Render2Real provide additional close precedents
for human-source motion and visual augmentation. ViKi is therefore positioned
as an investigation of a particular markerless multiview RGB-D route and its
operating limits. The intended contribution is to implement and characterise
that route across declared capture, embodiment, and rendering conditions,
including failures at their interfaces. The camera setup and target robot set
are experimental variables, rather than grounds for an unsupported claim of
first integration.

\section{Conclusion of the Review}
\label{sec:lit-conclusion}

The surveyed literature supports the proposed direction while making its
assumptions and evidence requirements explicit. The following implications
connect the review to the methodology:

\begin{enumerate}[leftmargin=2em]
  \item \emph{Retain a reusable human representation.} Articulated hand
        observations and hand--object relationships can support different
        target abstractions. A wrist pose and scalar aperture are suitable
        for some grippers, while dexterous hands require richer transfer.
  \item \emph{Compare retargeting formulations under matched conditions.}
        Evaluate hand-only and interaction-aware objectives using the same
        recordings and robot placements. Report tracking error, motion
        regularity, and constraint violations separately.
  \item \emph{Treat geometry and appearance as connected representations.}
        Specify the static background, dynamic objects, human removal, and
        replacement robot, together with their metric alignment and visibility
        handling. Photorealistic appearance does not establish contact validity.
  \item \emph{Measure capture quality on the relevant setup.} Calibration,
        depth validity, camera placement, and timing affect both perception
        and novel-view synthesis. External reference data and local rig
        measurements answer different questions.
  \item \emph{Validate the interfaces between stages.} Alongside independent
        stage metrics, check coordinate frames, time alignment, object state,
        and agreement between rendered robot motion and exported actions.
  \item \emph{Characterise accepted expansion rather than nominal size.}
        Record configuration coverage, rejected cases, and preparation and
        generation effort. Keep those findings distinct from physical task
        success and policy-learning performance.
\end{enumerate}

These choices define the proposed dataset-generation study without presupposing
its outcome. Chapter~\ref{ch:method} develops the relevant formulations, while
implementation and experimental evidence determine which capabilities can
ultimately be claimed.
